\documentclass[11pt,a4paper]{article}
\usepackage[T1]{fontenc}
\usepackage{lmodern}
\usepackage[margin=27mm]{geometry}
\usepackage{amsmath,amssymb,amsthm}
\usepackage{booktabs}
\usepackage{microtype}
\usepackage{enumitem}
\usepackage[hidelinks]{hyperref}
\hypersetup{pdftitle={Congruences for compositional power series},pdfauthor={Dev Goyal},pdfsubject={OEIS A396099 and related formal power series; coefficient congruences}}
\setlength{\parindent}{1.2em}
\setlength{\parskip}{3pt}
\setlength{\emergencystretch}{2em}
\newtheorem{theorem}{Theorem}[section]
\newtheorem{lemma}[theorem]{Lemma}
\newtheorem{corollary}[theorem]{Corollary}
\theoremstyle{remark}
\newtheorem{remark}[theorem]{Remark}
\newcommand{\Z}{\mathbb Z}
\newcommand{\N}{\mathbb N}
\newcommand{\coeff}[2]{[x^{#1}]#2}
\title{Congruences for compositional power series}
\author{Dev Goyal}
\date{24 October 2026}
\begin{document}
\maketitle

\begin{center}
\fbox{\begin{minipage}{0.88\linewidth}
\textbf{Contribution and review status.}
ChatGPT (OpenAI) developed the mathematical arguments and the accompanying
code autonomously in response to prompts from Dev Goyal. The author
initiated the investigation and requested checks, extensions, and
formalization. Independent human review of this new work is pending.
Dev Goyal is the sole human author. He would bear responsibility for
any final release; this draft awaits his review. The formalization
package records its own checked coverage;
the written proofs below do not assume that every statement has been
verified by a proof assistant.
\end{minipage}}
\end{center}

\begin{abstract}
For each positive integer $r$, let $F_r(x)=x+O(x^2)$ be the formal power
series determined by
$F_r(x)=x+F_r^{\circ r}(x)F_r^{\circ(r+1)}(x)$, where the superscript
$\circ k$ denotes composition $k$ times. We determine every iterate
modulo $2r$. For odd $r$ the reduction of $F_r$ is $x/(1-x)$; for even
$r$ it is $x/(1-x)+rx^4/(1-x)^3$. This proves the four congruence
assertions recorded for OEIS A396099. It also proves seven iterate
assertions for A396798, corrects the starting index of a residue block
in that entry, and proves the two assertions for A396807. The
$r=3$ case was proved previously by Sundaram. For a second family,
defined by $G_q(x-xG_q(qx)/q)=x$, we prove a uniform congruence modulo
$q+1$ and an exact parity criterion. These give six assertions
recorded in A393856--A393859. The arguments use
coefficientwise uniqueness, an elementary substitution identity, and
rational comparison series. All required formal-series facts are
proved below.
\end{abstract}

\section{The question and the main result}

The sequence OEIS A396099 begins
\[
1,\ 1,\ 5,\ 39,\ 383,\ 4357,\ 55057,\ 754223,\ldots.
\]
Paul D.\ Hanna defined its generating function $A(x)=\sum_{n\ge1}a_nx^n$
by
\begin{equation}\label{eq:a396099}
A(x)=x+A(A(x))\,A(A(A(x))).
\end{equation}
The large coefficients have a simple apparent pattern: they are odd,
and their remainders modulo four repeat $1,3,3,1$ starting with the
third coefficient. The entry also records two assertions involving
compositions of $A$ \cite{oeis099}.

We write $U^{\circ k}$ for a compositional iterate, so that
$U^{\circ2}(x)=U(U(x))$. Ordinary powers, such as $U(x)^2$, mean
multiplication. Our three basic identities are
\begin{align}
A(x)&\equiv\frac{x}{1-x}+\frac{2x^4}{(1-x)^3}\pmod4,
\label{eq:mainA}\\
A^{\circ2}(x)&\equiv x+2x^2\pmod4,
\label{eq:mainB}\\
A(x)-xA^{\circ3}(x)&\equiv x+\frac{2x^3}{1-x}\pmod4.
\label{eq:mainC}
\end{align}
A congruence of power series means that the corresponding coefficients
have the same remainder. Thus \eqref{eq:mainB} says that every
coefficient of the second iterate above degree two is divisible by
four. Equation \eqref{eq:mainC} says that every coefficient of the
displayed difference above degree two has remainder two. The initial
$x$ in this identity is part of the full series statement.

The explanation of \eqref{eq:mainA} is equally explicit. Since
\[
\frac{x^4}{(1-x)^3}
=x^4+3x^5+6x^6+10x^7+15x^8+\cdots,
\]
we obtain
\begin{equation}\label{eq:a-coeff}
a_n\equiv1+2\binom{n-2}{2}\pmod4\qquad(n\ge2),
\end{equation}
where $\binom02=\binom12=0$. The binomial coefficient is odd precisely
at $n\equiv0,1\pmod4$ in this range. The residues of $a_n$, from
$n=3$ onward, are therefore $1,3,3,1$ repeated. Every positive-index
coefficient is odd.

These facts arise from a whole family. For $r\ge1$, let $F_r$ satisfy
\begin{equation}\label{eq:family}
F_r=x+F_r^{\circ r}F_r^{\circ(r+1)},\qquad F_r=x+O(x^2).
\end{equation}
Existence and uniqueness are proved in Section~\ref{sec:formal}.
Theorem~\ref{thm:general} gives a formula for all nonnegative iterates
of $F_r$ modulo $2r$. Section~\ref{sec:integer} extends it to integer
iterates by using the compositional inverse.

\paragraph{Attribution and status.}
Hanna recorded the instances considered here and their conjectures
\cite{oeis099,oeis797,oeis798,oeis807}. Sundaram proved the
$r=3$ case, including all integer iterates, in September 2026
\cite{sundaram}. His comparison with $x/(1-x)$ and uniqueness argument
are directly relevant prior work. The odd-$r$ argument below extends
that calculation. The even-$r$ case requires the extra correction
$rx^4/(1-x)^3$. The entries for $r=2,4,5$ still labelled the statements
discussed here as conjectures when checked on 24 October 2026.
This describes the checked sources; it does not establish exclusive
historical priority for the results.

\section{Formal power series and uniqueness}\label{sec:formal}

\subsection{Definitions}

For a commutative ring $R$ with identity, a formal power series is an
expression $U(x)=\sum_{n\ge0}u_nx^n$ with coefficients in $R$.
The notation $R[[x]]$ denotes the set of these series. Equality means
equality of every coefficient. No numerical value is assigned to $x$,
so analytic convergence is not involved.

Addition is coefficientwise, and multiplication is defined by
\[
\coeff{n}{(UV)}=\sum_{j=0}^n u_jv_{n-j}.
\]
When $V(0)=0$, substitution is defined by
$U(V(x))=\sum_{j\ge0}u_jV(x)^j$. Every coefficient receives only
finitely many contributions. Associativity and the distributive laws
then follow coefficient by coefficient from polynomial identities.

A series is \emph{normalized} if its constant coefficient is zero
and its coefficient of $x$ is one. Thus it has the form $x+O(x^2)$.
The notation $O(x^N)$ denotes a series with zero coefficients in
degrees below $N$. In particular, $U-V=O(x^N)$ says that $U$ and $V$
agree in those degrees. Products satisfy
$O(x^a)O(x^b)=O(x^{a+b})$.

For a normalized $U$, define $U^{\circ0}=x$ and
$U^{\circ(k+1)}=U\circ U^{\circ k}$ for $k\ge0$.
Each iterate is normalized. If $U-V=O(x^N)$, then
\begin{equation}\label{eq:iterate-congruence}
U^{\circ k}-V^{\circ k}=O(x^N)\qquad(k\ge0).
\end{equation}
Indeed, this follows by induction on $k$: split a difference of
compositions into the change in the outer series and the change in
the inner series. Substitution by a normalized series preserves the
degree of vanishing, and factoring $P^j-Q^j$ shows that a change of
order $N$ in the inner argument creates no term of smaller degree.

Reduction of integer coefficients modulo an integer commutes with
these operations, because each coefficient is a finite sum of
products. A series with constant coefficient one has a multiplicative
inverse: if $W(0)=0$, then
\[
(1+W)^{-1}=\sum_{j\ge0}(-W)^j.
\]
This is well defined coefficientwise, and successive terms cancel
upon multiplication. It justifies every fraction used below, whose
denominator has constant coefficient one.

\subsection{The canonical solution}

\begin{lemma}\label{lem:unique}
For every commutative ring $R$ with identity and every integer $r\ge1$,
there is exactly one normalized series $F\in R[[x]]$ satisfying
\[
F=x+F^{\circ r}F^{\circ(r+1)}.
\]
\end{lemma}

\begin{proof}
For normalized $U$, set
$\Phi_r(U)=x+U^{\circ r}U^{\circ(r+1)}$.
This series is normalized. Suppose $U-V=O(x^N)$, with $N\ge2$.
By \eqref{eq:iterate-congruence}, corresponding iterates differ by
$O(x^N)$. Each iterate starts with $x$, so
\begin{align*}
\Phi_r(U)-\Phi_r(V)
={}&(U^{\circ r}-V^{\circ r})U^{\circ(r+1)}\\
&+V^{\circ r}(U^{\circ(r+1)}-V^{\circ(r+1)})
=O(x^{N+1}).
\end{align*}
Thus applying $\Phi_r$ improves agreement by one degree.

For existence, start with $U_0=x$ and define
$U_{j+1}=\Phi_r(U_j)$. We have $U_1-U_0=O(x^2)$; the preceding
calculation gives $U_{j+1}-U_j=O(x^{j+2})$ by induction. Hence every
coefficient eventually stops changing. Define $F$ by those stabilized
coefficients. Each coefficient of $\Phi_r(U)$ depends on only finitely
many coefficients of $U$, so the limit satisfies $F=\Phi_r(F)$.

For uniqueness, two normalized solutions agree below degree two.
Whenever they agree below degree $N$, the displayed calculation and
their defining equations force agreement below degree $N+1$.
Induction gives agreement in every degree.
\end{proof}

The proof uses no division, so it applies to residue rings even when
the modulus is composite. In particular it constructs the integer
series $F_r$ in \eqref{eq:family}. Its image under any ring
homomorphism is the canonical solution over the target ring, by
uniqueness. We can therefore identify its reduction modulo $2r$ by
finding any normalized solution of the reduced equation.

The first coefficients are useful for checking the indexing:
\begin{equation}\label{eq:firstcoeff}
F_r(x)=x+x^2+(2r+1)x^3+O(x^4).
\end{equation}
The product on the right side of \eqref{eq:family} starts with $x^2$,
so the second coefficient is one. It follows by induction that
$F_r^{\circ j}=x+jx^2+O(x^3)$. Multiplying the iterates for $j=r$
and $j=r+1$ gives coefficient $2r+1$ in degree three.

\section{Comparison series and a substitution identity}

Define, for an integer $c$,
\begin{equation}\label{eq:tH}
T_c(x)=\frac{x}{1-cx},\qquad T=T_1,\qquad
H(x)=\frac{x^4}{(1-x)^3}.
\end{equation}
These are integer power series. Direct substitution gives
\begin{equation}\label{eq:t-group}
T_c\circ T_d=T_{c+d}.
\end{equation}
Consequently $T^{\circ k}=T_k$ for every nonnegative integer $k$.
The coefficient of $x^n$ in $T_k$ is $k^{n-1}$ for $n\ge2$,
and its coefficient of $x$ is one.

The formal derivative is defined algebraically by
$U'(x)=\sum_{n\ge1}nu_nx^{n-1}$. It satisfies the product and
substitution rules, which can be proved first for monomials and then
coefficientwise. We need the following rule over a ring in which a
particular element has square zero.

\begin{lemma}\label{lem:taylor}
Let $\varepsilon\in R$ satisfy $\varepsilon^2=0$, and let
$u,v\in xR[[x]]$. For every $S\in R[[x]]$,
\[
S(u+\varepsilon v)=S(u)+\varepsilon vS'(u).
\]
\end{lemma}

\begin{proof}
For every nonnegative integer $j$, the binomial theorem gives
$(u+\varepsilon v)^j=u^j+j\varepsilon vu^{j-1}$, with the second
term interpreted as zero for $j=0$. Terms containing two or more
factors of $\varepsilon$ vanish. Multiply by the coefficient of
$x^j$ in $S$ and sum. All substitutions are well defined because
their inner arguments have zero constant coefficient.
\end{proof}

There is no division by a factorial in this rule. It therefore works
in composite-modulus rings.

\begin{lemma}\label{lem:identities}
The following identities hold in $\Z[[x]]$:
\begin{align}
H(x)T'(T(x))+H(T(x))&=\frac{2x^4}{(1-2x)^3},
\label{eq:identity1}\\
(1-x)H(x)-x^2T(x)-x^3T'(x)&=-\frac{2x^3}{1-x}.
\label{eq:identity2}
\end{align}
\end{lemma}

\begin{proof}
We have $T'(x)=(1-x)^{-2}$, so
\[
H(x)T'(T(x))=\frac{x^4}{(1-x)(1-2x)^2},\qquad
H(T(x))=\frac{x^4}{(1-x)(1-2x)^3}.
\]
Their sum has numerator $x^4((1-2x)+1)=2x^4(1-x)$ over the
common denominator $(1-x)(1-2x)^3$. This proves
\eqref{eq:identity1}. For the second identity, use denominator
$(1-x)^2$ to obtain the numerator
\[
x^4-x^3(1-x)-x^3=-2x^3(1-x).
\]
Cancellation gives \eqref{eq:identity2}.
\end{proof}

Finally, expanding three geometric series shows that
\begin{equation}\label{eq:Hcoeff}
\coeff{n}{H}=
\begin{cases}
0,&n<4,\\
\binom{n-2}{2},&n\ge4.
\end{cases}
\end{equation}
The coefficient of $x^d$ in $(1-x)^{-3}$ counts triples of
nonnegative integers with sum $d$. For each first entry $i$ there
are $d-i+1$ choices for the other two, giving
$(d+1)(d+2)/2=\binom{d+2}{2}$. This proves
\eqref{eq:Hcoeff} without an analytic expansion.

\section{The general congruence}

\begin{theorem}\label{thm:general}
Let $r\ge1$ and $k\ge0$ be integers. Then
\begin{equation}\label{eq:general-congruence}
F_r^{\circ k}(x)\equiv
\begin{cases}
\displaystyle\frac{x}{1-kx},&r\text{ odd or }k\text{ even},\\[7pt]
\displaystyle\frac{x}{1-kx}+\frac{rx^4}{(1-x)^3},
&r\text{ even and }k\text{ odd},
\end{cases}
\pmod{2r}.
\end{equation}
\end{theorem}

\begin{proof}
First suppose $r$ is odd. Since $r+1$ is even,
\[
(1-rx)(1-(r+1)x)
=1-(2r+1)x+r(r+1)x^2\equiv1-x\pmod{2r}.
\]
Using \eqref{eq:t-group}, we obtain
\[
x+T^{\circ r}T^{\circ(r+1)}
\equiv x+\frac{x^2}{1-x}=T\pmod{2r}.
\]
Thus $T$ is the normalized solution of the reduced equation.
Lemma~\ref{lem:unique} gives $F_r\equiv T\pmod{2r}$, and
composition yields $F_r^{\circ k}\equiv T_k$ for all $k\ge0$.

Now suppose $r$ is even. Work in the coefficient ring
$R=\Z/(2r)\Z$, and set $f=T+rH$. The image of $r$ in this ring
has square zero, because $r^2$ is a multiple of $2r$.
Lemma~\ref{lem:taylor} and \eqref{eq:identity1} give
\begin{align*}
f^{\circ2}
&=T(T+rH)+rH(T+rH)\\
&=T^{\circ2}+r\bigl(HT'(T)+H(T)\bigr)\\
&=T^{\circ2}+\frac{2rx^4}{(1-2x)^3}
=T^{\circ2}\quad\text{in }R[[x]].
\end{align*}
Because $r$ is even, repeated composition of this identity gives
\[
f^{\circ r}=T^{\circ r}=\frac{x}{1-rx}=x+rx^2
\quad\text{in }R[[x]].
\]
The last equality again uses $r^2=0$ in $R$. Hence
\[
f^{\circ(r+1)}=f(x+rx^2)=f+rx^2f'.
\]
We can now verify the defining equation. After expanding and discarding
terms with a factor of $r^2$, its residual is
\begin{align*}
f-\bigl(x+f^{\circ r}f^{\circ(r+1)}\bigr)
&=(1-x)f-x-rx^2f-rx^3f'\\
&=r\bigl((1-x)H-x^2T-x^3T'\bigr)\\
&=-\frac{2rx^3}{1-x}=0\quad\text{in }R[[x]],
\end{align*}
where \eqref{eq:identity2} proves the penultimate equality.
By uniqueness, the reduction of $F_r$ is $f$.

For $k=2j$, the equality $f^{\circ2}=T^{\circ2}$ gives
$f^{\circ k}=T^{\circ k}=T_k$. For $k=2j+1$, it gives
\[
f^{\circ k}=f\circ T_{2j}=T_k+rH(T_{2j}).
\]
The integer series $T_{2j}$ is congruent to $x$ modulo two.
Therefore $H(T_{2j})-H(x)$ has even integer coefficients, and
multiplication by $r$ makes it zero modulo $2r$. Thus
$f^{\circ k}=T_k+rH$ in $R[[x]]$, as required.
\end{proof}

\begin{corollary}\label{cor:coeff}
Write $h_n=\coeff{n}{H}$ as in \eqref{eq:Hcoeff}, and let
$\epsilon(r,k)$ equal one if $r$ is even and $k$ is odd, and zero
otherwise. For every $r\ge1$, $k\ge0$, and $n\ge2$,
\[
\coeff{n}{F_r^{\circ k}}
\equiv k^{n-1}+r\epsilon(r,k)h_n\pmod{2r}.
\]
The coefficient of degree zero is zero, and that of degree one is one.
Every positive-index coefficient of $F_r$ is odd.
\end{corollary}

\begin{proof}
The coefficient formula follows from the geometric series for $T_k$
and \eqref{eq:Hcoeff}. At $k=1$, the correction is absent for odd
$r$ and is a multiple of the even integer $r$ for even $r$.
Since $2$ divides $2r$, all positive-index coefficients are odd.
\end{proof}

\section{The four assertions for A396099}

Take $r=2$, so $A=F_2$. The cases $k=1$ and $k=2$ of
Theorem~\ref{thm:general} immediately give \eqref{eq:mainA} and
\eqref{eq:mainB}. The parity assertion follows from
Corollary~\ref{cor:coeff}.

For completeness, the phase of the residue block can be checked
directly. For $n\ge2$, put $b_n=\binom{n-2}{2}$, with value zero
when $n=2,3$. Then
\[
b_{n+4}-b_n=4n-2,
\]
which is even. The initial values at $n=2,3,4,5$ are $0,0,1,3$.
Thus $b_n$ is even at $n\equiv2,3\pmod4$ and odd at
$n\equiv0,1\pmod4$. Equation~\eqref{eq:a-coeff} gives
\[
a_{4j+3}\equiv1,\quad a_{4j+4}\equiv3,\quad
a_{4j+5}\equiv3,\quad a_{4j+6}\equiv1\pmod4
\qquad(j\ge0).
\]
This establishes the recorded residue cycle with its correct start.

It remains to prove \eqref{eq:mainC}. In $\Z/4\Z[[x]]$, write
$f=T+2H$. The proof of Theorem~\ref{thm:general} gave
$f^{\circ2}=x+2x^2$, so Lemma~\ref{lem:taylor} gives
\[
f^{\circ3}=f+2x^2f'.
\]
Consequently
\begin{align*}
f-xf^{\circ3}
&=(1-x)f-2x^3f'\\
&=x+2\bigl((1-x)H-x^3T'\bigr)\\
&=x+2\frac{x^4-x^3}{(1-x)^2}
=x-\frac{2x^3}{1-x}
=x+\frac{2x^3}{1-x}\pmod4.
\end{align*}
Since $A\equiv f\pmod4$, this proves \eqref{eq:mainC}. Its
coefficients above degree two are all congruent to two. All four
assertions in A396099 have now been proved.

\section{Integer iterates and exact composition order}\label{sec:integer}

The statements for A396797 and A396807 include negative iterate
indices. We give the extension explicitly.

\begin{lemma}\label{lem:inverse}
Every normalized series over a commutative ring has a unique
normalized two-sided inverse under composition.
\end{lemma}

\begin{proof}
For $U=x+\sum_{n\ge2}u_nx^n$, seek
$V=x+\sum_{n\ge2}v_nx^n$ with $U(V)=x$.
The coefficient equation in degree $n$ has the form $v_n$ plus
an expression involving only earlier coefficients of $V$.
Indeed, in $V^j$ for $j\ge2$, a coefficient of degree $n$ uses
only indices at most $n-j+1<n$. Thus every $v_n$ is determined
without division, giving a unique right inverse.
Apply the construction to $V$ to obtain $W$ with $V\circ W=x$.
Associativity gives
$U=U\circ(V\circ W)=(U\circ V)\circ W=W$.
It follows that $V\circ U=x$ too.
\end{proof}

For an integer $k<0$, define $F_r^{\circ k}$ by composing the inverse
$-k$ times. The usual law of integer powers holds in this composition
group. Reducing coefficients commutes with the inverse, because it
preserves both inverse identities and the normalized inverse is unique.

\begin{corollary}\label{cor:order}
The reduction of $F_r$ modulo $2r$ has exact order $2r$ under
composition. Formula~\eqref{eq:general-congruence} remains valid
for every integer $k$, including negative values. In particular,
\[
F_r^{\circ k}\equiv x\pmod{2r}
\quad\Longleftrightarrow\quad 2r\mid k.
\]
\end{corollary}

\begin{proof}
For nonnegative $k$, the coefficient of $x^2$ in the formula of
Theorem~\ref{thm:general} is $k$ modulo $2r$, because $H$ has
no terms below degree four. Equality with $x$ therefore forces
$2r\mid k$. Conversely, if $2r\mid k$, then $k$ is even,
$T_k\equiv x\pmod{2r}$, and the correction term is absent.
This proves the order statement.

For any integer $k$, choose the unique $j$ with
$0\le j<2r$ and $j\equiv k\pmod{2r}$. In the composition group
modulo $2r$, the order statement gives
$F_r^{\circ k}=F_r^{\circ j}$. Also $T_k\equiv T_j\pmod{2r}$,
and $k$ and $j$ have the same parity, because $2r$ is even.
The formula for $j$ is therefore exactly the required formula for
$k$. The same coefficient-of-$x^2$ argument proves the final
equivalence for all integer $k$.
\end{proof}

\section{Related OEIS assertions}

\subsection{A396798: seven iterate assertions and a phase correction}

The entry A396798 is the case $r=4$ \cite{oeis798}. Let $B=F_4$.
Theorem~\ref{thm:general} gives
\[
B^{\circ k}\equiv T_k+4\,\mathbf1_{\{k\text{ odd}\}}H\pmod8.
\]
For even $k=2,4,6,8$, the geometric series truncates modulo eight
as shown in the table below. For odd $k$, we have $k^2\equiv1\pmod8$.
Using the parity pattern of $h_n$, the coefficients from degree two
repeat $k,1,k+4,5$ modulo eight. These calculations prove all seven
iterate assertions recorded in the entry.

\begin{center}
\begin{tabular}{@{}cl@{}}
\toprule
Iterate & Reduction modulo eight\\
\midrule
$B^{\circ2}$ & $x+2x^2+4x^3$\\
$B^{\circ3}$ & coefficients $3,1,7,5$ repeated from degree two\\
$B^{\circ4}$ & $x+4x^2$\\
$B^{\circ5}$ & coefficients $5,1,1,5$ repeated from degree two\\
$B^{\circ6}$ & $x+6x^2+4x^3$\\
$B^{\circ7}$ & coefficients $7,1,3,5$ repeated from degree two\\
$B^{\circ8}$ & $x$\\
\bottomrule
\end{tabular}
\end{center}

At $k=1$, the correct coefficient block is $1,1,5,5$ beginning
at \emph{degree two}. The entry begins this block at degree one;
that phase is false. Equation~\eqref{eq:firstcoeff} gives
\[
B=x+x^2+9x^3+O(x^4).
\]
Thus its third coefficient has remainder one modulo eight, whereas
the stated phase would require five. The corrected assertion is
\[
\coeff{4j+2}{B}\equiv1,\quad
\coeff{4j+3}{B}\equiv1,\quad
\coeff{4j+4}{B}\equiv5,\quad
\coeff{4j+5}{B}\equiv5\pmod8
\qquad(j\ge0).
\]
It follows for all indices from the general theorem.

\subsection{A396807 and the previously proved A396797 case}

For $r=5$, the series is OEIS A396807 \cite{oeis807}. The odd-parameter
case of the theorem, together with Corollary~\ref{cor:order}, gives
\[
F_5^{\circ k}\equiv\frac{x}{1-kx}\pmod{10}\qquad(k\in\Z).
\]
Its coefficient of $x$ is one. For $n\ge2$, its coefficient of
$x^n$ is $k^{n-1}$ modulo ten. Taking $k=1$ gives remainder one
for every positive-index coefficient of $F_5$. These are the two
assertions in the entry. Stating degree one separately also handles
$k=0$ without relying on an unexplained expression $0^0$.

The same odd-parameter argument at $r=3$ recovers Sundaram's result
for A396797 \cite{oeis797,sundaram}. We include it as an instance
of the general formula, with the prior result explicitly credited.

\section{A scaled inverse family}\label{sec:scaled}

A second family shows how the same uniqueness principle can lead to
a different arithmetic pattern. For a positive integer $q$, consider
the normalized series $G_q$ satisfying
\begin{equation}\label{eq:scaled-source}
G_q\left(x-\frac{xG_q(qx)}q\right)=x.
\end{equation}
Hanna's entries A393856, A393857, A393858, and A393859 correspond
respectively to $q=4,5,6,7$
\cite{oeis856,oeis857,oeis858,oeis859}. Each entry proposes that
every positive-index coefficient has remainder one modulo $q+1$.
The entries for $q=4,6$ also propose that a coefficient is odd
exactly when its index is a power of two. These six statements
remained labelled conjectures in the sources checked on
24 October 2026.

\subsection{Removing the apparent division}

If $U(x)=\sum_{n\ge1}u_nx^n$, define
\[
L_q(U)(x)=\sum_{n\ge1}q^{n-1}u_nx^n.
\]
For a nonzero integer $q$, this is exactly $U(qx)/q$. The coefficient
definition makes sense over every commutative ring, without requiring
that $q$ be invertible. Equation~\eqref{eq:scaled-source} therefore
means
\begin{equation}\label{eq:scaled-nodiv}
G_q\bigl(x-xL_q(G_q)(x)\bigr)=x.
\end{equation}

\begin{lemma}\label{lem:scaled-unique}
For any commutative ring $R$ with identity and any $q\in R$, equation
\eqref{eq:scaled-nodiv} has a unique normalized solution in $R[[x]]$.
Its construction commutes with ring homomorphisms.
\end{lemma}

\begin{proof}
Write $G=x+\sum_{n\ge2}g_nx^n$ and
$V=x-xL_q(G)$. The series $V$ is normalized. The coefficient $g_n$
inside $V$ first occurs in degree $n+1$, because of the initial
factor $x$. Thus the coefficient of degree $n$ in $V$ depends
only on $g_2,\ldots,g_{n-1}$.

In $G(V)=V+\sum_{j\ge2}g_jV^j$, the term $g_nV^n$ contributes
exactly $g_n$ in degree $n$. The terms with $j<n$ depend only on
earlier coefficients, and those with $j>n$ contribute nothing in
that degree. Therefore the equation in degree $n$ has the form
$g_n+P_n=0$, with $P_n$ determined by earlier coefficients and $q$.
This gives existence and uniqueness by recursion, without division.
It also shows that applying a ring homomorphism to the coefficients
preserves the equation; uniqueness identifies the resulting solution.
\end{proof}

\subsection{The congruence and parity theorems}

\begin{theorem}\label{thm:scaled-congruence}
For every positive integer $q$,
\[
G_q(x)\equiv\frac{x}{1-x}\pmod{q+1}.
\]
In particular, every coefficient of positive degree is congruent to
one modulo $q+1$.
\end{theorem}

\begin{proof}
Work in $R=\Z/(q+1)\Z$, where $q=-1$. For $T=x/(1-x)$,
\[
L_{-1}(T)=x-x^2+x^3-\cdots=\frac{x}{1+x}.
\]
Hence
\[
x-xL_{-1}(T)=x-\frac{x^2}{1+x}=\frac{x}{1+x},
\]
and direct substitution gives
\[
T\left(\frac{x}{1+x}\right)
=\frac{x/(1+x)}{1-x/(1+x)}=x.
\]
Thus $T$ is the normalized solution of the reduced equation.
Lemma~\ref{lem:scaled-unique} identifies it with the reduction of
$G_q$. This argument also covers composite moduli such as six
and eight.
\end{proof}

\begin{theorem}\label{thm:scaled-parity}
If $q$ is a positive even integer, then
\[
G_q(x)\equiv\sum_{j\ge0}x^{2^j}\pmod2.
\]
Equivalently, its coefficient of degree $n\ge0$ is odd if and only
if $n$ is a power of two. If $q$ is odd, all positive-index
coefficients of $G_q$ are odd.
\end{theorem}

\begin{proof}
For even $q$, reduction modulo two gives $L_q(U)=x$ for every
normalized $U$: all factors $q^{n-1}$ with $n\ge2$ vanish.
Therefore the reduced equation is
\[
G_q(x+x^2)=x.
\]
Let $S(x)=\sum_{j\ge0}x^{2^j}$ over $\Z/2\Z$.
The strictly increasing exponents make this a well-defined formal
series. Repeated squaring in characteristic two gives
$(x+x^2)^{2^j}=x^{2^j}+x^{2^{j+1}}$, and hence
\[
S(x+x^2)
=\sum_{j\ge0}\bigl(x^{2^j}+x^{2^{j+1}}\bigr)=x.
\]
Every term beyond the first occurs twice and cancels. By uniqueness,
the reduction of $G_q$ is $S$. Its support is exactly the powers
of two, including $1=2^0$ and excluding zero.

If $q$ is odd, then $q+1$ is even, so
Theorem~\ref{thm:scaled-congruence} already gives oddness of every
positive-index coefficient.
\end{proof}

The four congruence applications and two parity applications are
summarized below. They all concern the solution of the original
equation \eqref{eq:scaled-source}, not a separately defined sequence
chosen to fit the remainders.

\begin{center}
\begin{tabular}{@{}cll@{}}
\toprule
$q$ & OEIS entry & Consequences\\
\midrule
$4$ & A393856 & remainder one modulo five; power-of-two parity\\
$5$ & A393857 & remainder one modulo six\\
$6$ & A393858 & remainder one modulo seven; power-of-two parity\\
$7$ & A393859 & remainder one modulo eight\\
\bottomrule
\end{tabular}
\end{center}

\section*{Verification and responsibility}

The proofs in this article apply to all indices. Exact finite
computations in the accompanying package reconstruct coefficients
from the defining equation and check the reductions, compositions,
and displayed examples. Such computations can expose indexing or
implementation errors; they are separate from the all-index proofs.

The accompanying Lean development proves the four original A396099
assertions for the canonical integer series. It also proves the
general congruence for every positive parameter and every natural or
integer iterate, including the exact composition order; the seven
iterate assertions for A396798, its corrected coefficient phase and
a counterexample to the stated phase; both A396807 assertions,
including negative iterates; and all six assertions for
A393856--A393859. The definitions use actual formal power series,
and negative iterates use their compositional inverses.

The final package's coverage table, explicit theorem types, fresh
build records, and transitive axiom audit are the authoritative
record of formal verification. The Lean proofs sometimes take a
different algebraic route from the exposition: in particular,
exact ring identities replace the derivative calculation in the
even-parameter argument. This does not claim a separate formalization
of every intermediate lemma in the article. Finite coefficient
checks remain corroborating computations, not all-index proofs.
Historical priority, source status, and authorship are outside
Lean's certification.

ChatGPT (OpenAI) developed the mathematical arguments and the
accompanying computational and formalization work autonomously from
the author's prompts. The author initiated the project and requested
verification, extensions, and revisions; the author did not supply
the proof strategy or intermediate mathematical arguments.
Independent human review of this new work remains pending. Dev Goyal
is the sole human author. The intended responsibility for a final
release rests with him, including its mathematical claims,
references, attribution, and any errors; this draft awaits his
review. The verification status of this work is independent of any
earlier project by the author.

\begin{thebibliography}{9}
\bibitem{oeis099}
P. D. Hanna, \emph{A396099}, The On-Line Encyclopedia of Integer
Sequences. \url{https://oeis.org/A396099}. Accessed 24 October 2026.

\bibitem{oeis797}
P. D. Hanna, \emph{A396797}, The On-Line Encyclopedia of Integer
Sequences. \url{https://oeis.org/A396797}. Accessed 24 October 2026.

\bibitem{oeis798}
P. D. Hanna, \emph{A396798}, The On-Line Encyclopedia of Integer
Sequences. \url{https://oeis.org/A396798}. Accessed 24 October 2026.

\bibitem{oeis807}
P. D. Hanna, \emph{A396807}, The On-Line Encyclopedia of Integer
Sequences. \url{https://oeis.org/A396807}. Accessed 24 October 2026.

\bibitem{sundaram}
Shiva Priyan Shanmuga Sundaram, \emph{A Proof of the Congruence Conjectures
for OEIS A396797}, September 2026.
\url{https://oeis.org/A396797/a396797.pdf}.
Accessed 24 October 2026.
\bibitem{oeis856}
P. D. Hanna, \emph{A393856}, The On-Line Encyclopedia of Integer
Sequences. \url{https://oeis.org/A393856}. Accessed 24 October 2026.

\bibitem{oeis857}
P. D. Hanna, \emph{A393857}, The On-Line Encyclopedia of Integer
Sequences. \url{https://oeis.org/A393857}. Accessed 24 October 2026.

\bibitem{oeis858}
P. D. Hanna, \emph{A393858}, The On-Line Encyclopedia of Integer
Sequences. \url{https://oeis.org/A393858}. Accessed 24 October 2026.

\bibitem{oeis859}
P. D. Hanna, \emph{A393859}, The On-Line Encyclopedia of Integer
Sequences. \url{https://oeis.org/A393859}. Accessed 24 October 2026.
\end{thebibliography}
\end{document}
